%% Draft v0.7 — 2026-06-13 (review pass 5, post-referee-feedback).
%% Changes from v0.6:
%%   - sec 3.1 free-alpha Einasto resolves the fixed-shape inflation (54%->32%,
%%     three-backbone agreement); table note b rewritten; abstract de-claims
%%     "robustly require" -> robust-to-cap-and-backbone, threshold-dependent.
%%   - sec 3.2 tie-aware W/L/T scoring (2-BIC band): 36/0/5 Burkert, 27/7/7 NFW,
%%     32/7/2 Einasto; N=1-vs-Burkert degeneracy flagged N/A (tie artifact
%%     FIXED in pipeline, not noted in text).
%%   - sec 3.7 rewritten: power-control quantified; the two warped cases adjudicated
%%     on warp-aware curves (NGC 2841 reproduces dBIC=11; NGC 3521 honest
%%     non-reproduction dBIC=-7; NGC 1003 power-limited, quantified); NGC 3741
%%     supporting reproduction added. "instrument-limited" bucket 3->1.
%%   - Abstract: inner-component disclaimer promoted (outer is the diagnostic
%%     signal; no claim on inner displacement); "domain" defined as radial mass
%%     component, no 3D shape implied.
%%   - Refs: Di Teodoro & Fraternali 2015 (now cited in sec 3.7 as the
%%     forward-modeling tool), Walter et al. 2008 (cited in sec 3.7).
%%   - Sample-definition and exclusion-rule TODOs closed.
%%   - v0.7 FIGURES: three figures inserted (Fig 1 worked example+shell sec3.1;
%%     Fig 2 concordance sec3.4; Fig 3 forced-fit null sec3.5). Source-order numbering.
%%   - AASTeX 7 port: documentclass aastex701; author block uses v7 syntax
%%     (\author[orcid] then \affiliation then \email).
%%   - CLERICAL FILLED 2026-06-24: contact email, acknowledgments + SPARC citation,
%%     multi-tracer cross-ref (-> in preparation); "real" softened to
%%     "supported by convergent evidence" in conclusions. REMAINING: Zenodo DOI
%%     (data-availability block, from the git->Zenodo release).
%% Target: ApJL (AASTeX v7.0.1). Letter format.
%% Numbers frozen from halo-spheres/ analysis of 2026-06-05; sec 3.7-3.8 from
%% program2/ analysis of 2026-06-12; see README.md and program2/SCOPING.md.

\documentclass[twocolumn]{aastex701}

\newcommand{\vflat}{V_{\rm flat}}
\newcommand{\reff}{R_{\rm eff}}
\newcommand{\dbic}{\Delta{\rm BIC}}
\providecommand{\doi}[1]{\href{https://doi.org/#1}{doi:#1}}
\newcommand{\msun}{M_\odot}

\shorttitle{Residual radial structure in SPARC halos}
\shortauthors{Bibb}

\begin{document}

\title{Method-Independent Detection of Residual Radial Structure in SPARC
Dark Matter Halos Beyond Single Smooth Equilibrium Profiles}

\author[0009-0004-1153-2464]{Ron Bibb}
\affiliation{Independent Researcher, Lilburn, GA, USA}
\email[show]{ronbibb@gmail.com}

\begin{abstract}
Single smooth equilibrium halo profiles leave statistically significant,
radially organized residuals in roughly one third of SPARC rotation curves.
We introduce a minimum-order radial-domain decomposition: the dark-matter
contribution is approximated by a minimum-order sum of tiered cored support
components (Burkert form), with the order selected by the Bayesian
Information Criterion. Across 121 galaxies, a second radial domain is required
in $\sim$35\% of cases --- a fraction robust to confining component scales to
the measured radii (a constraint that simultaneously collapses all
three-domain solutions) and stable across three halo families (Burkert 35\%,
NFW 37\%, free-shape Einasto 32\%), though, as for any information-criterion
threshold, the precise fraction is threshold-dependent. The central result is
convergence: galaxies requiring a second domain coincide with galaxies
independently flagged by a localized Gaussian-shell decomposition (odds ratio
38; mass-adjusted 37.5), and the domain crossover radius tracks the
independently fitted shell radius ($\rho = 0.73$, permutation $p = 0.004$,
surviving half-light-radius control). Forced two-domain fits to single-domain
galaxies collapse to degenerate configurations, providing an internal null.
Where the structure is strong enough to test, the multi-domain selection
reproduces on rotation curves re-derived from the original radio cubes by
independent software --- decisively ($\dbic = 121$ for NGC~2403) and at the
same crossover radius --- establishing that it is a property of the galaxy,
not of the fitting basis. In all 41 two-domain systems the single-profile
scale radius lies between the two domain scales (median factor 2.2): single
smooth fits are systematically biased summaries, not merely incomplete ones.
We interpret the second component as a method- and data-independent marker of
a second cored radial mass domain (the term ``domain'' denoting a radial mass
component, with no three-dimensional shape implied, which rotation curves
cannot constrain). Of the two components, the diagnostic signal is carried by
the \emph{outer}: it remains consistent with single-domain halo scaling,
whereas forced decompositions of single-domain galaxies cannot reproduce that
behavior. We make no claim about the inner component's relation to the halo
scaling locus --- a displacement forcing reproduces --- and remain agnostic
as to physical origin.
\end{abstract}

\keywords{Dark matter (353) --- Galaxy rotation curves (619) --- Galaxy dark
matter halos (1880) --- Astrostatistics (1882)}

%% ------------------------------------------------------------------
\section{Introduction} \label{sec:intro}

Rotation curves of late-type galaxies are conventionally decomposed into
baryonic contributions plus a single smooth equilibrium dark matter halo,
described by cuspy profiles motivated by $\Lambda$CDM simulations
\citep[NFW;][]{Navarro1996}, observationally motivated cored profiles
\citep[Burkert;][]{Burkert1995}, or intermediate forms
\citep[Einasto;][]{Einasto1965}. The adequacy of a \emph{single} smooth
profile is usually assessed only through global fit quality. Localized
departures --- coherent residual features at particular radii --- have been
discussed largely in connection with baryonic features
\citep[``Renzo's rule'';][]{Sancisi2004}, but a systematic, model-agnostic
census of whether single smooth halos suffice is lacking.

In a companion analysis (Bibb 2026, Zenodo, \doi{10.5281/zenodo.20263016}) we
modeled localized residual structure in SPARC \citep{Lelli2016} rotation
curves with Gaussian mass shells superposed on a Burkert backbone, finding
BIC-selected shells in roughly half the sample. That decomposition is used
here only as a structurally independent comparison basis, not as a source of
physical claims. The present Letter asks a deliberately minimal,
complementary question: \emph{what is the minimum number of smooth,
concentric, cored radial domains required by each rotation curve?} Because
any sufficiently flexible component will absorb residuals --- so that
inferred ``structure'' may reflect the chosen basis rather than the galaxy
--- the credibility of either decomposition rests on validation. We therefore
test the answer against the structurally unrelated shell decomposition,
against mass-driven confounds, against an internal forced-fit null, and ---
most stringently --- against rotation curves of the same galaxies derived
independently from the original radio data cubes.

Four properties make the result reportable. First, it is robust: the
multi-domain fraction is unchanged when components are forbidden from
extrapolating beyond the measured radii, while unstable three-domain
solutions are eliminated by the same constraint. Second, it is
backbone-independent: NFW and Burkert nulls give the same fraction. Third, it
is method-independent: two unrelated decompositions flag the same galaxies and
localize structure to consistent radii, and this concordance is not explained
by galaxy mass. Fourth, and most stringently, it is \emph{data}-independent
where it can be tested: for galaxies whose rotation curves can be re-derived
from the original radio cubes by independent software, the multi-domain
selection reproduces --- decisively, and at the same radius. A fit improvement
on one curve, however constrained the basis, cannot predict the outcome of an
independent derivation from a different instrument; that this prediction
succeeds is the strongest available evidence that the structure is a property
of the galaxy rather than of the parameterization. We emphasize equally what
the analysis does \emph{not} find, and we make no claim about the physical
origin or three-dimensional geometry of the detected structure.

%% ------------------------------------------------------------------
\section{Data and Methods} \label{sec:methods}

\subsection{Sample and baryonic model}

We use the SPARC database \citep{Lelli2016}, restricted to galaxies with
quality flag $Q \le 2$ and at least six rotation-curve points after removing
points interior to the first reliably measured radius: 123 galaxies meet
these criteria, of which 121 yield converged fits. Baryonic contributions
adopt fixed mass-to-light ratios $\Upsilon_{\rm disk} = 0.5$ and
$\Upsilon_{\rm bul} = 0.7$ at $3.6\,\mu$m with the gas contribution at unity,
combined with signed squares so that $V^2_{\rm bar} = V_{\rm gas}|V_{\rm
gas}| + 0.5\,V_{\rm disk}|V_{\rm disk}| + 0.7\,V_{\rm bul}|V_{\rm bul}|$.
Velocity uncertainties are floored at $1\,{\rm km\,s^{-1}}$. Points at which
the baryonic model exceeds the observed velocity ($V^2_{\rm obs} \le V^2_{\rm
bar}$) are excluded from the dark-matter fit, matching the companion analysis.

\subsection{Minimum-order concentric-domain decomposition}
\label{sec:tiers}

The inferred dark-matter contribution is approximated by a minimum-order sum
of cored radial support components of Burkert form,
\begin{equation}
V^2_{\rm DM}(r) = \sum_{i=1}^{N} V^2_{\rm B}(r;\,\rho_{0,i},\,a_i),
\end{equation}
subject to ordering constraints that enforce a tiered structure and break
exchange degeneracies: $\rho_{0,i}$ strictly decreasing and $a_i$ strictly
increasing outward. A soft prior ties the total enclosed dark mass at the
outermost measured point to the dynamically inferred value. Fits use bounded
nonlinear least squares with 20--24 random restarts; model order is selected
by ${\rm BIC} = \chi^2 + 2N \ln n_{\rm pts}$, evaluated for $N = 1$--$4$.
At $N = 1$ the model reduces exactly to a single Burkert halo, providing the
null hypothesis within the same machinery.

Two variants are run. In the \emph{unconstrained} variant, scale radii are
free. In the \emph{capped} variant, every component is required to satisfy
$a_i \in [r_{\rm min}, r_{\rm max}]$, the measured radial range, so that no
component's characteristic scale lies where there are no data. The cap is not
a technical variant but a \emph{stability discriminator}: structure that
survives it is resolved by the data, while structure that collapses is
identified as extrapolation artifact. It thereby provides an internal
falsification mechanism for each candidate domain.

\subsection{Comparison methods}

\emph{Gaussian-shell decomposition.} The companion method models localized
residuals as Gaussian mass shells of center $r_s$, width $\sigma$, and mass
$M_s$ on a single Burkert backbone, with shell count BIC-selected and a width
constraint $\sigma/r \le 0.4$ that defines locality. The preferred fractional
width is cap-invariant: the median $\sigma/r \simeq 0.22$ is unchanged as the
allowed ceiling is raised from 0.4 to 1.0, with only 4\% of shells near the
ceiling at a cap of 1.0. The shell and domain decompositions share data but
use structurally unrelated basis functions --- a localized annular perturbation
versus a global cored profile --- and are fit independently.

\emph{Backbone generalization.} The minimum-order analysis of
\S\ref{sec:tiers} is repeated with NFW and Einasto in place of Burkert. For
Einasto we run two variants: a fixed-shape null ($\alpha = 0.16$) and a
free-shape variant in which $\alpha$ is an additional shared parameter
($k = 2N + 1$, with the extra parameter charged to the BIC), bounded to the
physical range $\alpha \in [0.1, 0.3]$.

%% ------------------------------------------------------------------
\section{Results} \label{sec:results}

\subsection{A robust two-domain population} \label{sec:twodomain}

Table~\ref{tab:order} summarizes the BIC-selected model order.
Unconstrained: 78/38/5 galaxies select $N = 1/2/3$ (64/31/4\%). Capped:
79/41/0 (one galaxy selects $N = 4$). The multi-domain fraction is
essentially unchanged, 36\% versus 35\%, and 95\% (41/43) of multi-domain
galaxies retain $N \ge 2$ under the cap. By contrast, \emph{all five}
three-domain solutions collapse: in every case the third component was an
inner spike below the first measured point or an outer ramp beyond the last,
i.e.\ extrapolation rather than measurement. The data support at most two
resolved radial domains.

The decision margins are strongly asymmetric. Where two domains win, they win
decisively: median $\dbic = +43.3$ over the single profile, with 23/42 cases
exceeding $+30$ and a maximum of $+1115$. Where the single profile wins, it
never wins decisively: the median losing margin is $+4.8$, 81\% of losses are
within $\dbic < 6$, and \emph{no} galaxy prefers a single Burkert by
$\dbic > 10$. The single-domain population is therefore adequate rather than
demanded, and the quoted 35\% should be read as threshold-dependent in the
usual information-criterion sense.

A direct consequence concerns what a single profile reports when two domains
are present. In all 41 capped two-domain systems, the single-Burkert scale
radius lies between the inner and outer domain scales, with a median
$a_{\rm Burkert}/a_{\rm in} = 2.19$, indicating that a single smooth
backbone compresses the two-domain structure into a biased intermediate scale
rather than recovering the inner domain plus a residual correction. Smooth
fits can therefore not only conceal radial structure but systematically bias
inferred core scales, by a factor of $\sim$2 in this sample. Figure~%
\ref{fig:worked} shows a representative example.

\begin{figure}
\centering
\includegraphics[width=\linewidth]{figures/fig1_composite.png}
\caption{Representative two-domain decomposition (NGC~2403, SPARC). The observed
rotation curve (points) and baryonic contribution (dotted) are shown with the best
single-Burkert fit (dashed) and the BIC-selected two-sphere composite (thick solid);
the inner and outer cored components are shown individually (thin solid). The single
profile cannot reproduce both the inner rise and the outer profile simultaneously,
whereas the two-domain composite does ($\chi^2_\nu$ improving from $\simeq 19$ to
$\simeq 3$). The domain crossover (vertical line, $\simeq 2.3$~kpc) lies close to the
radius of the independently fitted Gaussian shell (marked, $\simeq 1.0$~kpc) for this
galaxy, previewing the population-level concordance of Figure~\ref{fig:concordance}.
This is also one of the systems for which the selection reproduces on a rotation curve
re-derived from the original radio cubes by independent software (\S\ref{sec:dataindep}).
\label{fig:worked}}
\end{figure}

\begin{deluxetable}{lcccc}
\tablecaption{BIC-selected number of radial domains \label{tab:order}}
\tablehead{\colhead{Variant} & \colhead{$N{=}1$} & \colhead{$N{=}2$} &
\colhead{$N{=}3$} & \colhead{$N{\ge}2$ frac.}}
\startdata
Unconstrained & 78 & 38 & 5 & 36\% \\
Capped ($a \in [r_{\rm min}, r_{\rm max}]$) & 79 & 41 & 0\tablenotemark{a} & 35\% \\
NFW backbone (unconstrained) & 76 & 36 & 9 & 37\% \\
Einasto backbone\tablenotemark{b} & 56 & 47 & 18 & 54\% \\
\enddata
\tablenotetext{a}{One galaxy selects $N = 4$.}
\tablenotetext{b}{Fixed $\alpha = 0.16$. The inflated fraction is a
fixed-shape artifact: freeing $\alpha$ (physical range $[0.1, 0.3]$, with the
extra parameter BIC-charged) returns 32\% (\S\ref{sec:backbone}), consistent
with Burkert and NFW. The fitted $\alpha$ rails at the upper bound,
indicating the shape freedom is spent absorbing inner structure with one
steep component rather than adding a domain --- the mechanism of the
fixed-shape inflation. Even an unphysically flexible single Einasto
($\alpha$ to 0.5) reaches only 28\%, below the Burkert fraction, so the
two-domain detection is conservative against an over-flexible null.}
\end{deluxetable}

\subsection{Backbone independence and comparison to single profiles}
\label{sec:backbone}

With an NFW backbone the multi-domain fraction is 37\%, statistically
indistinguishable from the Burkert 35--36\%: the residual structure is not an
artifact of fitting a cored profile to cuspy halos or vice versa. The
fixed-shape Einasto null returns 54\%, but this is a rigidity artifact:
freeing the Einasto shape parameter $\alpha$ returns 32\% (Table~\ref{tab:order},
note b). The two-domain fraction is therefore stable at 32--37\% across all
three halo families once each is given a fair shape, with the fixed-$\alpha$
Einasto 54\% identified as the outlier rather than evidence of extra structure.

For the capped two-domain galaxies, the composite outperforms every
single-profile alternative on BIC, including the parameter penalty. Scoring
wins/losses/ties with a 2-BIC significance band (a margin below which the
comparison is non-discriminating), the composite records 36/0/5 against a
single Burkert, 27/7/7 against a single NFW, and 32/7/2 against a single
Einasto: it never loses to a single Burkert and loses to NFW or Einasto in at
most seven cases. (For single-domain galaxies the composite is by
construction identical to a single Burkert, so that comparison is degenerate;
it is reported as non-applicable rather than scored.) The residual structure
is not absorbable by exchanging one smooth profile family for another.

\subsection{Galaxy-level concordance with the shell method, controlling for
mass} \label{sec:concordance}

The central result is cross-method agreement. Of the 42 capped multi-domain
galaxies, 39 are independently shell-bearing; of the 79 single-domain
galaxies, 59 are shell-free (Fisher exact $p < 10^{-4}$; odds ratio 38;
overall agreement 81\%).

Because both structure indicators correlate with galaxy mass, we test
whether the concordance is a mass artifact. In logistic regression for
multi-domain selection with $\log M_\star$ as the base predictor, shell-%
bearing status adds likelihood-ratio $\chi^2 = 49.7$ ($p < 10^{-4}$) ---
an order of magnitude beyond any other covariate --- while $T$-type,
effective surface brightness, $R_{\rm HI}$, gas fraction, and bulge presence
(LR $\chi^2 = 0.01$, $p = 0.93$) are entirely absorbed by mass, and
$\vflat$ contributes only weakly ($p = 0.013$). The bulge null additionally
disarms the concern that the second component merely compensates for
bulge--disk mass-to-light errors: if it did, bulge presence would predict
multi-domain selection at fixed mass, and it does not. Stratifying into $\log M_\star$ tertiles, the shell--domain
association holds within every bin (Fisher $p = 0.011$, $<10^{-4}$,
$<10^{-4}$), and the Cochran--Mantel--Haenszel mass-adjusted odds ratio is
37.5 --- indistinguishable from the raw 38. At fixed mass, shell-bearing
galaxies are $\sim$37 times more likely to require a second radial domain.

\subsection{Radial-level concordance} \label{sec:radial}

Co-occurrence in the same galaxies could still mean the two methods describe
different features. We therefore compare \emph{where} they place structure.
For each capped two-domain galaxy we compute the crossover radius
$r_\times$ at which the outer component's velocity contribution overtakes the
inner, and compare it to the independently fitted shell radius (nearest
shell; $n = 28$ galaxies host both).

The two radii are strongly correlated: Spearman $\rho = 0.73$
($p < 10^{-4}$), with a median fractional offset
$|r_\times - r_s|/r_{\rm max} = 0.188$ against $0.387$ for permuted shell
assignments (permutation $p = 0.0038$). Both quantities partly track the
optical half-light radius ($\rho = 0.57$ for $r_\times$ vs.\ $\reff$), so we
control for it: residualizing both radii on $\reff$, the shell--crossover
correlation persists at partial $\rho = 0.62$ ($p = 4 \times 10^{-4}$). The
two methods localize structure to common radii for reasons beyond a shared
preference for the mid-disk. Notably, the capped fits show \emph{stronger}
radial concordance than the unconstrained fits ($\rho = 0.73$ vs.\ $0.68$;
permutation $p = 0.004$ vs.\ $0.047$), as expected if the cap removes
extrapolated components that scatter the crossover radius. The galaxy-by-galaxy
correspondence is shown in Figure~\ref{fig:concordance}.

\begin{figure}
\centering
\includegraphics[width=\linewidth]{figures/fig2_concordance.png}
\caption{Radial concordance between the two decompositions. For the $n=28$ galaxies
hosting both a BIC-selected second domain and an independently fitted shell, the domain
crossover radius $r_\times$ is plotted against the nearest shell radius $r_s$
(logarithmic axes); the solid line is $r_\times = r_s$. The two radii are strongly
correlated (Spearman $\rho = 0.73$), and the correlation survives control for the
optical half-light radius (partial $\rho = 0.62$), so the two methods localize structure
to common radii for reasons beyond a shared preference for the mid-disk.
\label{fig:concordance}}
\end{figure}

\subsection{An internal null: forced two-domain fits} \label{sec:forced}

As an internal control, we force $N = 2$ fits on the 79 galaxies whose
BIC-selected solution is single-domain. In 71/79 the added component yields
negligible likelihood improvement ($\Delta\chi^2 < 2$); only 8/79 improve
genuinely but fail the BIC penalty. The forced components collapse toward
degenerate configurations: median $a_{\rm out}/a_{\rm in} = 1.16$, versus
9.09 in selected two-domain systems ($p < 10^{-4}$), with the added
component's density suppressed by $\sim$4 dex. Moreover, the inner--outer
scale coupling present in selected systems survives control for galaxy size
(partial $\rho = 0.56$, $p = 10^{-4}$), whereas in forced fits it vanishes
(partial $\rho = 0.13$, $p = 0.25$). The decomposition therefore possesses an
identifiable null morphology, distinct from the selected multi-domain
population.

A further property of the forced fits sharpens this null into a positive
discriminant. Placing both the forced and the genuinely selected components
against the single-domain halo $\rho_0$--scale relation --- built locally in
$V_{\rm flat}$ to remove the mass-regime dependence --- the two component tiers
behave asymmetrically under forcing. The \emph{inner} components of forced and
selected fits are statistically indistinguishable in their displacement below the
relation (Mann--Whitney $p = 0.10$): forcing any second component drives an inner
term below the relation, so inner displacement alone is \emph{not} diagnostic of
genuine structure, and we make no claim on it. The \emph{outer} components, by
contrast, separate cleanly: in selected two-domain systems the outer component
remains comparatively close to the single-domain locus (median $-0.95\sigma$),
whereas in forced decompositions --- excluding degenerate fits that rail against
parameter bounds --- it is driven far below it (median $-3.44\sigma$;
Mann--Whitney $p = 2 \times 10^{-4}$, KS $p = 2 \times 10^{-4}$). Were the
decomposition mere fitting freedom, both tiers would respond to forcing alike;
instead one tier survives the forcing and the other does not. That asymmetry is
itself evidence that the outer domain of a selected system tracks a genuine halo
component rather than an artifact of adding parameters. Figure~\ref{fig:forced}
contrasts the selected and forced populations directly.

\begin{figure}
\centering
\includegraphics[width=\linewidth]{figures/fig3_forced_null.png}
\caption{The forced-fit null and its asymmetry. For single-domain galaxies, a second
component is \emph{forced} and compared with the components of genuinely BIC-selected
two-domain systems. The inner components are statistically indistinguishable in their
displacement below the single-domain $\rho_0$--scale relation (Mann--Whitney
$p = 0.10$): forcing reproduces the inner offset, so it is not diagnostic. The outer
components separate cleanly --- selected systems remain near the locus
(median $-0.95\sigma$) while forced fits are driven far below it (median $-3.44\sigma$;
$p = 2\times10^{-4}$). One tier survives forcing and the other does not, identifying the
outer domain as the component that tracks a genuine halo feature.
\label{fig:forced}}
\end{figure}

\subsection{What does not predict the second domain} \label{sec:negatives}

Three deliberately reported negative results bound the interpretation.

\emph{No privileged scale ratio.} The capped outer/inner scale ratio has
median $a_{\rm out}/a_{\rm in} \simeq 9$, but it does not correlate with the
dimensionless structural ratios $R_{\rm HI}/R_{\rm d}$ or
$r_{\rm max}/R_{\rm d}$ ($|\rho| \le 0.13$, $p > 0.4$); it correlates only
with absolute size and mass. The absolute scales track the measured extent
($a_{\rm out}$ vs.\ $r_{\rm max}$: $\rho = 0.80$). The decomposition scales
with the galaxy; it does not identify a special radius ratio.

\emph{No common surface-density manifold.} Pooled components follow a
$\rho_0$--$a$ anticorrelation of slope $-1.29$, superficially resembling the
constant surface-density relation \citep{Donato2009}, but inner and outer
components occupy significantly different surface-density levels
($\log \rho_0 a = 8.59$ vs.\ $7.98$; $p < 0.001$). The two tiers are not a
single scaling law sampled twice.

\emph{No gas-fraction signal at fixed mass.} Within mass tertiles,
multi-domain and single-domain galaxies have indistinguishable gas fractions
($p = 0.13$--$0.96$).

%% ------------------------------------------------------------------
\subsection{Reproduction on independently derived rotation curves} \label{sec:dataindep}

The preceding tests establish that the second domain is independent of the
fitting basis. They do not, by themselves, exclude the possibility that the
structure originates in the specific rotation curves published by SPARC --- for
example in a localized systematic of one velocity-field extraction. We therefore
test the strongest form of independence available: whether the multi-domain
selection reproduces when the rotation curve is re-derived from the original
radio data cubes by independent software, holding the baryonic model fixed so
that only the measured velocities change.

For the dwarf control galaxies DDO~154 and DDO~168, asymmetric-drift-corrected
circular-velocity curves derived by an independent group from the same survey
cubes (3D forward-modeling; \citealt{Iorio2017}) reproduce the single-domain
selection ($N=1$) in both cases. On the multi-domain side, we re-derived curves
for NGC~2403 and NGC~3198 directly from public THINGS-class cubes using
independent tilted-ring forward-modeling software
\citep[\textsuperscript{3D}Barolo;][]{DiTeodoro2015}, with geometry fixed to
literature tilted-ring solutions. Both reproduce the multi-domain selection
decisively: NGC~2403 selects $N=2$ at $\dbic = 121$ with the domain crossover at
$1.9$~kpc (SPARC-based value $\sim 2.3$~kpc), and NGC~3198 selects $N=2$ at
$\dbic = 55$. For NGC~2403 the independent curve also derives from different
observations than those underlying the SPARC curve, strengthening the test. A
third multi-domain galaxy, NGC~3741, reproduces at lower significance on an
independent VLA-ANGST extraction \citep[different telescope from the
WSRT-based SPARC curve;][]{Walter2008}: the multi-domain geometry recovers,
though the modest signal-to-noise of this dwarf limits the decision margin.

Two further multi-domain galaxies, NGC~2841 and NGC~3521, are strongly warped,
so a blind fixed-geometry extraction cannot test them; we instead use published
warp-aware tilted-ring rotation curves (DiskFit; Sellwood \& Spekkens, private
communication) derived from the same THINGS cubes. \textbf{NGC~2841
reproduces}: the warp-aware curve selects $N=2$ at $\dbic = 11$. Its earlier
failure on a blind extraction is attributable to the warp suppressing the
feature, which the warp-aware curve recovers. \textbf{NGC~3521 does not
reproduce}: its warp-aware curve prefers a single domain ($\dbic = -7$ against
two), an honest non-reproduction that we report as such. An injection--recovery
power control on the warp-aware curves confirms both verdicts are informative:
the two-domain feature these galaxies would need to show perturbs the curve by
$\sim$20--50~km\,s$^{-1}$, well above the $\sim$13--19~km\,s$^{-1}$ amplitude
their sampling and uncertainties can detect, so neither is power-limited.

A single galaxy, NGC~1003, is genuinely instrument-limited and we report it as
such with a quantified justification: its independent (HALOGAS) extraction
yields only 14 usable rings, and an injection--recovery power control shows that
the two-domain feature it would carry (peak deviation $\sim$3.3~km\,s$^{-1}$
over a single-domain fit) lies below the $\sim$7~km\,s$^{-1}$ amplitude its
sampling and noise can detect. Its non-reproduction is therefore uninformative
--- the curve is incapable of resolving the feature --- rather than a negative
result.

The adjudicated tally is thus four reproductions (NGC~2403, NGC~3198,
NGC~2841, NGC~3741) spanning three classes of independence (same
cubes/different software, different observations, different telescope), one
honest non-reproduction (NGC~3521), one power-limited indeterminate case
(NGC~1003, quantified), and two single-domain controls reproduced
(DDO~154/168). Reproduction holds in four of the five adjudicable galaxies;
the one failure and the one exclusion are reported explicitly rather than
absorbed into an undifferentiated ``instrument-limited'' category.

%% ------------------------------------------------------------------
\subsection{Relation to the shell decomposition: absorption by the domain model} \label{sec:absorption}

The galaxy- and radius-level concordance of \S\ref{sec:concordance}--\ref{sec:radial}
shows that the shell and domain methods flag the same galaxies at the same radii.
A sharper question is whether the localized shells carry information \emph{beyond}
the broad two-domain structure, or whether they are largely a localized
re-parameterization of it. We test this directly by freezing the BIC-selected
sphere-domain model of each galaxy and asking whether Gaussian shells, fit to the
residual of that frozen model using the identical shell machinery and BIC penalty
of the companion analysis, survive.

They mostly do not. Among the capped two-domain galaxies, 21 of 28 shells that
are selected against a single-Burkert backbone are absorbed once the second
domain is modeled --- they no longer improve the fit beyond the BIC penalty.
Population-wide, shell detections fall from 33 to 11 once the domain model is in
place. Only three galaxies retain a shell at $\dbic > 10$ against the frozen
two-domain backbone; in these, the surviving shell sits well away from the domain
crossover, consistent with a genuinely localized feature distinct from the broad
structure. The implication is twofold: most localized ``shells'' in the companion
analysis are the broad two-domain structure expressed in a localized basis ---
which is why a localized basis drew the objection that it merely improves the
fit --- and the domain decomposition is the more economical description, with a
small residual minority of genuinely localized features isolated for separate
study.

%% ------------------------------------------------------------------
\section{Discussion} \label{sec:discussion}

The evidential chain is: (1) a single smooth equilibrium profile ---
whether NFW or Burkert --- leaves BIC-significant residual structure in
$\sim$35\% of SPARC galaxies; (2) the structure is limited to two resolvable
radial domains, with three-domain solutions exposed as extrapolation by the
data-range constraint; (3) an unrelated localized-shell decomposition flags
the same galaxies, and the association is untouched by mass control; (4) the
two methods place structure at consistent radii, beyond both a permutation
null and a generic mid-disk scale; (5) no other measured galaxy property
predicts the second domain; (6) where it can be adjudicated, the selection
reproduces on rotation curves re-derived from the original radio cubes by
independent software in four of five cases (including warp-aware curves), with
the single non-reproduction and one power-limited case reported explicitly;
and (7) when the broad two-domain model is fit first, it absorbs the large
majority of the localized shells, identifying them as the same structure in a
localized basis rather than independent features.

We emphasize what is and is not being claimed. The present analysis does not
imply that galaxies physically contain discrete Burkert subhalos, and it does
not claim decomposition \emph{uniqueness}. It claims decomposition
\emph{convergence}: structurally independent decomposition bases repeatedly
recover the same subset of galaxies as inconsistent with a single smooth
equilibrium support profile. We accordingly interpret the second component as
a \emph{marker of a second radial support domain}, not as a literal nested
halo: the negative results of \S\ref{sec:negatives} argue specifically
against reading the components as physical objects with preferred scale
ratios or a shared structural relation. A preliminary non-parametric density
inversion on a demonstration subset is consistent with this picture --- the
freely fit $\rho(r)$ shows low-order curvature rather than fragmenting into
features at arbitrary radii --- though we defer its quantitative development.

A second boundary on interpretation is geometric and absolute. A rotation curve
constrains only the enclosed-mass profile $M({<}r)$ through $V^2(r) = G
M({<}r)/r$; it carries no information about how the mass producing a given
$M({<}r)$ is arranged in angle or in the third dimension. A second radial domain,
a spherical shell, a planar ring, a stream, a warp-induced redistribution, and a
localized deposit are indistinguishable in $V(r)$ if they enclose the same mass
within each radius. We therefore make no claim about the three-dimensional
geometry of the detected structure: the components of our decomposition are
radial mass-profile features, and the labels ``domain'' and ``shell'' denote
parameterizations, not shapes. Discriminating a centrally concentrated from a
localized geometry requires observables that carry angular information ---
azimuthal structure in the velocity field or in resolved tracers --- and is
deferred to analyses of that kind. The present claim is confined to what
rotation curves can support: a reproducible, basis-independent feature in the
enclosed-mass profile.

The forced-fit control of \S\ref{sec:forced} makes the basis-flexibility
question empirical. Forced two-domain fits to single-domain galaxies collapse
toward degenerate same-scale, vanishing-density configurations with
negligible likelihood gain; BIC-selected two-domain systems instead exhibit
separated inner and outer scales with residual scale coupling at fixed galaxy
size. The detected structure is therefore not a generic consequence of
basis-function flexibility: the basis has a recognizable signature for
``nothing here,'' and the selected population does not show it.

Nor should single-domain galaxies be read as negative detections. The margin
asymmetry of \S\ref{sec:twodomain} --- decisive multi-domain wins against
uniformly marginal single-profile wins --- indicates heterogeneous structural
complexity across the population: some systems clearly demand additional
radial structure, while many are merely \emph{adequately} described by a
single profile rather than positively inconsistent with further structure.
Whether the detected structure's origin is baryonic (e.g.\ features
propagated from the gas disk; \citealt{Sancisi2004}), kinematic, or a genuine
property of the dark matter distribution is not determined here, and we
decline to speculate. The single external fact available is that, at fixed
mass, the second domain co-occurs with localized shell structure at odds
ratio $\sim$37.

We address directly the objection that a constrained multi-component fit might
simply be an elaborate means of minimizing $\chi^2$, reproducible by any
sufficiently flexible function such as a high-order polynomial. Five properties of
the result bear on it, none of which a generic flexible fit possesses. First, the
selection does not run away with complexity: under the data-range cap the
population settles at 65\% single-domain, 35\% two-domain, and zero three-domain
(\S\ref{sec:twodomain}), whereas a basis whose advantage were mere flexibility
would consume the sample toward ever-higher order. Second, the forced-fit null
(\S\ref{sec:forced}) shows the basis declines to manufacture a component where
none is selected --- forced two-domain fits collapse, rail, and degenerate rather
than improving smoothly; a polynomial does not refuse. Third, that same forced
control resolves asymmetrically (\S\ref{sec:forced}): forcing reproduces the
inner component's below-relation offset but cannot reproduce the outer
component's, so part of the decomposition demonstrably tracks something a flexible
fit cannot fake. Fourth, the cross-method concordance
(\S\ref{sec:concordance}--\ref{sec:radial}) is independent of the fitting
basis entirely: a structurally unrelated shell decomposition flags the same
galaxies (mass-adjusted odds ratio 37.5) at approximately the same radii (the
shell radius tracks the domain crossover at partial $\rho = 0.62$), which two
analyses fitting noise with different bases would not do. Fifth, the
data-independence test (\S\ref{sec:dataindep}) shows the selection reproduces on
rotation curves re-derived from the original radio cubes by independent software,
decisively and at the same radius. A $\chi^2$-minimizing re-parameterization has
none of these five properties. The one place the curve-fitting concern does bite
--- the inner component's displacement below the halo scaling relation, which
forcing reproduces --- we therefore decline to interpret; the surviving scaling
signal is the outer component's, not the inner's.

Two limitations bound the claims. The concordant sample is modest
($\sim$40 galaxies at the galaxy level; 28 at the radial level). The
basis-independence and data-independence tests are partial: the former shares the
data across methods, and the latter is currently established for the subset of
galaxies whose curves can be re-derived from public cubes at sufficient quality,
with one multi-domain system (NGC~1003) power-limited and one honest
non-reproduction (NGC~3521) reported explicitly (\S\ref{sec:dataindep}).
Confirmation against independent tracers of the same galaxies (in preparation)
is the natural next test, and is also the route to the geometric discrimination
that rotation curves cannot provide.

%% ------------------------------------------------------------------
\section{Conclusions} \label{sec:conclusions}

A minimum-order concentric-domain decomposition of 121 SPARC rotation curves
shows that $\sim$35\% of galaxies require a second radial domain beyond any
single smooth equilibrium profile (NFW, Burkert, or Einasto), a result robust
to confining all components within the measured radii --- a constraint that
simultaneously eliminates all three-domain solutions as artifacts. The
detection is method-independent: an unrelated Gaussian-shell decomposition
identifies the same galaxies (mass-adjusted odds ratio 37.5) and consistent
radii (partial $\rho = 0.62$ controlling for the half-light radius), and when
the domain model is fit first it absorbs the majority of those shells,
identifying them as the same structure in a localized basis. Where it can be
tested, the selection is also data-independent: it reproduces on rotation curves
re-derived from the original radio cubes by independent software, decisively and
at the same radius. The
component parameters do not organize on a privileged scale ratio or a common
surface-density relation, and no galaxy property beyond mass and shell
structure --- including bulge presence --- predicts the second domain. Where
two domains are present, the single-profile scale radius is a biased
compromise between the two domain scales (41/41 betweenness; median factor
2.2), so smooth fits both conceal the structure and distort the inferred
core scale. The residual radial structure concealed
by single smooth halo fits is therefore supported by convergent evidence,
bounded, reproducible on independent data, and --- as to its three-dimensional
geometry and physical origin --- deliberately left open, to be addressed with
observables that carry the angular information rotation curves lack.

\begin{acknowledgments}
We thank the SPARC team for making their database publicly available. This
work made use of the Spitzer Photometry and Accurate Rotation Curves (SPARC)
database \citep{Lelli2016}. This research was conducted independently, without
institutional or grant support.
\end{acknowledgments}

\software{numpy, scipy, pandas, astropy, matplotlib}

%% [TODO: Data availability — Zenodo DOI for the halo-spheres analysis
%% package (scripts, per-galaxy tables, figures), including the companion shell
%% analysis (doi:10.5281/zenodo.20263016) and the independent-derivation
%% reproductions of sec 3.7.]

\begin{thebibliography}{}
\bibitem[Burkert(1995)]{Burkert1995} Burkert, A. 1995, ApJL, 447, L25
\bibitem[Di Teodoro \& Fraternali(2015)]{DiTeodoro2015} Di Teodoro, E.~M., \& Fraternali, F. 2015, MNRAS, 451, 3021
\bibitem[Donato et al.(2009)]{Donato2009} Donato, F., Gentile, G., Salucci,
P., et al. 2009, MNRAS, 397, 1169
\bibitem[Einasto(1965)]{Einasto1965} Einasto, J. 1965, Trudy Astrofiz. Inst.
Alma-Ata, 5, 87
\bibitem[Iorio et al.(2017)]{Iorio2017} Iorio, G., Fraternali, F., Nipoti, C., et al. 2017, MNRAS, 466, 4159
\bibitem[Lelli et al.(2016)]{Lelli2016} Lelli, F., McGaugh, S.~S., \&
Schombert, J.~M. 2016, AJ, 152, 157
\bibitem[Navarro et al.(1996)]{Navarro1996} Navarro, J.~F., Frenk, C.~S., \&
White, S.~D.~M. 1996, ApJ, 462, 563
\bibitem[Sancisi(2004)]{Sancisi2004} Sancisi, R. 2004, in IAU Symp. 220,
Dark Matter in Galaxies, ed. S. Ryder et al. (San Francisco: ASP), 233
\bibitem[Walter et al.(2008)]{Walter2008} Walter, F., Brinks, E., de Blok, W.~J.~G., et al. 2008, AJ, 136, 2563
\end{thebibliography}

\end{document}
